\begin{paramdefinition}[Frames, coarsening and sizes for the plane lattice]\label{hyp:coarsened-plane-frames}\label{par:plane-rank}
\hypertarget{source-plane-rank}{}
In the plane construction let \(r\) be the plane dimension and \(K\) its codimension in dimension \(k=r+K\). Write \(a\) for the lower bound on the boundary degrees \(b_x\).
Let \(J\) be the number of frames, \(N\) the number of tops in a
cycle, and \(N_s\) an upper bound for the number of letters in a
segment. In the count use the lifted frame endpoints \(s_j,z_j\)
with \(s_j\) the first position of frame \(j\). Segment \(j\)
starts at \(s_{j-1}\) and ends at \(z_j\). Write \(Q\) for the
original bound on \(p^{z_j-x}\). The coarsening uses
\[
 A=\frac{a}{8Yk^2N_s},\qquad L=\lfloor\log_p A\rfloor,\qquad
 P_{j,x}=p^{L+s_{j-1}-x},\qquad H_j=p^{s_j-s_{j-1}}.
\]
Indeed the extra power at segment \(j\) is the largest power of
\(p\) at most \(\lfloor A/p^{z_j-s_{j-1}}\rfloor\).

Since \(A\ge Q\),
this equals \(p^{L-(z_j-s_{j-1})}\). In particular \(P_{j,x}\le A\),
\(P_{j+1,x}=P_{j,x}H_j\) on frame \(j\), and \(\prod_j H_j=q\).
Use the separation bound \(p^{L_1-1}\ge128 k^2N_sY^2\) and the
size bound
\begin{equation}\label{eq:plane-count-size}
  a\ge32pYk^2N_sQ^2(E_0+m_*+1).
\end{equation}
These are the separation and size inequalities for the plane model.
The argument below uses them with \(p\ge13\), \(r\ge2\),
\(Y\ge1\), \(N_s+r\le K\), \(a\le b_x\le 2a\), and the lifted
frame assumptions of Section~\ref{sec:plane-lifts}.
Put \(D_k(b)=\binom{b+k-1}{k-1}\).
\end{paramdefinition}
For the selected cycle, the separation and size comparisons in the proof of Proposition~\ref{prop:host-small-prime-test} use \(Y=x^{26}\), \(p^{L_1-1}\ge256k^2NY^2\) and \(\log Q=o(\log(a/b'))\). They hold uniformly over the admissible frames.


\begin{lemma}[Frame and digit degrees in the coarsened lattice]\label{lem:plane-digit-frame-lattice}
For the coarsened lattice of Definition~\ref{hyp:coarsened-plane-frames},
the late degrees satisfy
\[
 a'_x=b_x-E'_x-E_x>0.
\]

The digit and outgoing choices produce the same total degree
\(a'_x\) in the true coordinate at each top, with incoming
degree \(b_{j+1,x}\ge W_{j+1,x}\ge kN_s+1\).
The corresponding segment word families are independent over
the residue field in the fixed digit and degree choices.
\end{lemma}

\begin{proof}
\proofstep{The coarse degrees at the later tops.}
At each late top \(x\) of segment \(j\), meaning one after the
incoming frame, form the ordinary plane determinant of the last
\(r\) preceding even rows and give it the exponent
\[
 g_{j,x}=\left\lfloor\frac{b_u}{k^2Yp^{z_j-u}}\right\rfloor,
 \qquad u=\text{the earliest of those rows}.
\]
Write \(\mathcal D\) for the product over all these determinants
with their original radices \(p^{z_j-y}\).

At a top \(y\) its
true degree \(E_y\) is at most \(4N_s a/(k^2Y)\). A determinant
contributes at most \(g_{j,x}p^{z_j-y}\le 2a/(k^2Y)\), since
\(y\ge u\). Each top occurs in at most two segments. Use coarsened
true degrees \(a'_x=b_x-E'_x-E_x\), where \(E'_x\le E_0\) is the
prior cleaning budget. In particular
\[
 a'_x\ge a/2,
\]
since
\(4N_s/k^2\le 1/4\) follows from \(N_s+2r\le k\),
\(r\ge2\), and \(E_0\le a/32\).

At each top there is a unique occurrence after the incoming
frame. Choose a coordinate digit vector \(\ell_x\in[0,P_{j,x})^k\)
at this occurrence, with total \(|\ell_x|\equiv a'_x\pmod{P_{j,x}}\).
All digit choices have \(|\ell_x|<kP_{j,x}\le a'_x\). Write
\[
 n_x=(a'_x-|\ell_x|)/P_{j,x}.
\]
At a top strictly between frames
the degree in the segment is \(n_x\).

At an outgoing frame row
of segment \(j\) choose a degree \(e_x\in[0,rH_j)\) with
\(e_x\equiv n_x\pmod{H_j}\). Use degree \(b_{j+1,x}=(n_x-e_x)/H_j\)
when the same top reappears in the incoming frame. In fact
\[
 b_{j+1,x}>\frac{a'_x}{P_{j+1,x}}-k-r,
 \qquad
 W_{j+1,x}=\left\lfloor\frac{a'_x}{P_{j+1,x}}\right\rfloor-k-r-1,
 \qquad b_{j+1,x}\ge W_{j+1,x}.
\]
The strict inequality uses \(H_j\ge1\) and \(e_x<rH_j\). Our size
bounds give \(a'_x/P_{j+1,x}\ge4Yk^2N_s\), so
\[
 W_{j+1,x}\ge m_0=kN_s+1
\]

for \(m_0=kN_s+1\). The digit and frame degrees thus add to exactly
\(a'_x\) at every top, with nonnegative incoming padding.

\paragraph{The segment words.}
Fix the digits of the intervening tops in segment \(j\). The pre-outgoing path degree at \(x\) is
\[
 B_x=\begin{cases}n_x,&\text{intervening top},\\ W_{j,x}-m_0,&\text{incoming row}.\end{cases}
\]
Write \(S=\sum_x B_x\)
for this sum before the outgoing frame. Consider paths \(\mu\) in
\(k\) rows with these degrees. The last \(r\) steps are vector
steps of variable size. Just after the \(i\)-th outgoing step
require that the first \(i\) entries are all \(S\), with indices
starting at one.

Write \(\nu,\kappa\) for the restrictions to the
first \(r\) and the last \(K\) entries. Each is a signed Pieri
path with sizes determined by its weight changes. Each covector step creates at most one further negative entry. Thus
\[
 \#\{i:\mu_i<0\}\le N_s\le K-r,
\]
so the first \(r\) entries of the dominant weight \(\mu\) stay nonnegative. They stay at most \(S\) before the outgoing
frame by the sum of previous vector degrees. They stay at most
\(S\) afterward by the head condition and dominance. The
complement has zero degrees in the incoming frame because the
first \(r\) steps are vectors.

Impose the admissible path inequalities. At every
late top \(x\), \(\nu_r\) before the step is at least \(2B_u/3\)
for the earliest \(u\) of the preceding \(r\) vectors. The
complement increase at a vector and the plane decrease at a
covector are at most
\[
 \frac{\mathcal L a}{pkP_{j,x}}+
   \mathbf 1_{x\text{ even}}\frac{4a}{P_{j,u}},
 \qquad \mathcal L=C_{\mathcal L} N_s^2\log a.
\]
Let \(\mathcal T_j\) be the set of such paths with a choice of
standard terminal tableau for \(\kappa\). Its size depends only on
the intervening degrees and the fixed data.

For the plane word use
the tableau of its terminal rectangle \((S,\ldots,S)\), which is
unique. Let \(d_i\) be the full strip size at outgoing row \(i\).
It is at most \(2S\). At a vector step from \(\mu\) to \(\mu'\),
interlacing gives \(\sum_h(\mu'_h-\mu_h)\le\mu'_1-\mu_k\).
Here \(\mu'_1=S\), and \(\mu_k\ge-S\) since the total of all
covector degrees is at most \(S\).

\proofstep{Counting the outgoing decorations.}
For each member of \(\mathcal T_j\), choose outgoing decoration
exponents \(v_{i,h}\ge0\) for the \(r\) plane coordinates of
outgoing row \(i\), with
\begin{equation}\label{eq:plane-outgoing-decoration}
 v_{i,h}+2S<H_j,\qquad d_i+\sum_h v_{i,h}=e_i.
\end{equation}
Use the product of the plane and complement words, the outgoing
monomial with these exponents, and the same incoming factor for
all paths at the chosen degrees. The incoming factor is the
plane determinant of the incoming frame raised to \(m_0\), times
one coordinate of each incoming row to degree \(b_{j,x}-W_{j,x}\).
The product is in the indicated lattice \(M_j\). The signed strip
interlacing at the plane boundary supplies the two inequalities
between plane and complement at each late cut. The incoming determinant power shifts every plane coordinate of the weight by \(m_0\).

If the old minimum is \(\nu_r\), then
\[
 \min_i(\nu_i+m_0)=\nu_r+m_0.
\]
The other added vector factors can only increase this minimum.
The path bound on \(\nu_r\) before the step gives the absolute
lower cutoff independently. The normal degree of the product
at a vector is the complement strip size and at a covector it
is the plane strip size. The path inequalities give the normal
cutoff. Finally the exponents in each outgoing plane coordinate
are less than \(H_j\) by \eqref{eq:plane-outgoing-decoration} and
the bound
\[
 d_i\le2S.
\]


\proofstep{Independence in fixed degrees.}
These products are independent at each fixed set of degrees,
also over the residue field. First omit their common nonzero
incoming factor. In the \(i\)-th outgoing plane row specialize
all but coordinate \(r+1-i\) to zero. For each column
\(c\in[1,S]\) of the plane word, the height increases by zero
or one at each outgoing step. A rise wedges the sole surviving
coordinate there. Reading back from the final full exterior
subset, remove each of the coordinates at its rise. This leaves
precisely the terminal subset for the column at the pre-outgoing
cut, the complement of its rise coordinates. The coefficient is
\(1\) or \(-1\) by the wedge rule.

If \(R_c\) is the set of outgoing rises in column \(c\),
its pre-outgoing terminal subset is
\[
       [1,r]\setminus\{r+1-i:i\in R_c\}.
\]
Thus the specialization of the product word is, up to sign,
\[
 \prod_{i=1}^r x_{i,r+1-i}^{\,\#\{c:i\in R_c\}}
\]
times the
truncated plane word with those terminal subsets and the whole
complement word. The variables \(x_{i,r+1-i}\) are the surviving
outgoing coordinates.
The terminal subsets of the truncated plane word are standard.
At every outgoing cut the height on the left is at least the
height on the right.

If their current heights are \(h_L\ge h_R\), the remaining rises to the common terminal height obey
\[
 r-h_L\le r-h_R.
\]
Let \(D_L,D_R\) be the sets of removed coordinates in two adjacent columns. The surviving coordinates run in reverse order, so for every initial interval \([1,b]\), 
\[
 |D_L\cap[1,b]|\le|D_R\cap[1,b]|.
\]
 The retained
subsets therefore have weakly increasing entries across a row.
The truncated plane path and its subsets recover all
plane heights, since each missing coordinate specifies the time
of one rise. The two paths and their tableaux are therefore
recovered from the standard words in the variables other than
the outgoing plane variables.

To include the decorations in this argument, take the terms of
least total degree in the discarded coordinates in a putative
relation. For each tuple of exponents in those coordinates and
in the surviving outgoing plane coordinates, the displayed
specializations are independent by the standard word basis
in the remaining variables. Once the word indices are known, let \(a_i\) be their fixed antidiagonal exponents. The observed exponents \(b_i\) then recover each decoration by
\[
 d_i=b_i-a_i.
\]
This proves
independence and also explains why it holds with the same
coefficient ring as the residue rank calculation.

\end{proof}

\paragraph{Weight of paths at the segment cutoffs.}
\label{par:plane-typical}
\begin{lemma}[Weight of the paths obeying the cutoffs]
\label{lem:good-plane-path-count}
For segment \(j\) of Definition~\ref{hyp:coarsened-plane-frames}, let \(\mathcal T_j\) be the paths satisfying the degree cutoffs, with a choice of standard terminal tableau for the complement weight \(\kappa\). With \(D_k(b)=\binom{b+k-1}{k-1}\), their number satisfies
\begin{equation}\label{eq:plane-good-path-count}
 |\mathcal T_j|\ \ge\ (1-a^{-N_s^2})
    \prod_{x\text{ between frames}}D_k(n_x)
    \prod_{x\text{ incoming}}D_k(W_{j,x}-m_0).
\end{equation}
\end{lemma}

\begin{proof}
We prove this at the simultaneous cutoffs above. Let \(m\le N_s\)
be the number of letters in the segment and number the first
\(m-r\) steps before the outgoing frame. Use the hypotheses
\(k,\mathcal L\ge1088N_s^2\log a\). From the degree-loss bound
\(E'_x+E_x\le5b_x/(kY)\) we have \(a'_x\ge9b_x/10\).
For either kind of pre-outgoing row,
\[
 \frac{4a}{5P_{j,x}}\le B_x\le\frac{2a}{P_{j,x}},\qquad B_x\ge k.
\]
In fact the floor and padding lose at most
\[
 4kN_s\le b_x/(10P_{j,x}),
\]

since \(b_x/P_{j,x}\ge8Yk^2N_s\). Thus degrees of successive
rows of the same parity grow by at least
\[
 \rho=p^2/5>13.
\]

For any letter \(x\), each preceding pre-outgoing degree is at
most \(4\beta_x/p\), where \(\beta_x=a/P_{j,x}\), by strict
increase of the top positions.

First count all paths with the outgoing head conditions,
weighted by the number of tableaux of the terminal complement.
Before the outgoing frame the signed Pieri rules decompose the
tensor product of the indicated symmetric powers of the standard
module and its dual. Each path to a weight \(\lambda\) has weight
\(\dim V_\lambda\). After \(t\) outgoing steps remove the first \(t\)
entries, all equal to \(S\), and call the remaining row list the
tail.

At the next step its old and new tails \(\eta,\eta'\) obey
\[
 \eta_i\ge\eta'_i\ge\eta_{i+1}.
\]
Conversely every such
interlacing tail reinstates a vector strip on inserting the
extra head equal to \(S\), because all entries were at most
\(S\). Thus at the outgoing steps the ordinary branching rule
from \(\mathrm{GL}_{k-t}\) to \(\mathrm{GL}_{k-t-1}\) preserves the
total weight. These multiplicity-one rules and the Weyl dimension
formula apply to rational weights by a determinant twist and
duality. \cite{fulton-harris}
The last tail is the complement, whose tableau count is its
Weyl dimension. The total count before imposing the good path
inequalities is exactly
\[
 R=\prod_{x\text{ pre-outgoing}}D_k(B_x).
\]


\proofstep{The constraints on a path and its two signs.}
For a given path and sign \(\epsilon\in\{+,-\}\), let
\(\alpha^\epsilon\) be the nonincreasing list of positive parts
or of absolute negative parts of its tail. Up to the outgoing
frame the tail is the full list of \(k\) rows. Let \(\mathcal B^\epsilon\)
be the list of previously used pre-outgoing degrees of the
corresponding sign, with vectors positive and covectors negative.
For every \(T\ge0\) put
\[
 \delta_\epsilon(T)=\sum_{B\in\mathcal B^\epsilon}\min(B,T)
                    -\sum_h\min(\alpha^\epsilon_h,T).
\]
Each of these deficits starts at zero and is nondecreasing along
the path.

At a pre-outgoing step of the corresponding sign let
\(b\) be the strip size and \(\alpha,\alpha'\) the old and
new lists for this sign. Monotonicity and interlacing give
\[
 0\le\sum_h\bigl(\min(\alpha'_h,T)-\min(\alpha_h,T)\bigr)
           \le \min(b,T).
\]
The \(b\) bound sums the coordinate increments. For the
\(T\) bound use \(\alpha'_{h+1}\le\alpha_h\) and telescope. At a step of the opposite sign the
second sum can only decrease. During the outgoing frame both
lists of parts decrease entry by entry by the tail interlacing,
and the reference lists stay fixed.

Write
\[
 \Xi_x=\mathcal L\beta_x/(pk).
\]
We bound the weight of
paths violating any of the following conditions at a late top
\(x\): a covector decreases positive mass by at most \(\Xi_x\),
a vector decreases negative mass by at most \(\Xi_x\), and the
last plane entry before either step is at least \(2B_u/3\).
Here \(u\) is the earliest of the preceding \(r\) vectors, which
is before the outgoing frame. These conditions imply the good
path inequalities. At a covector all plane entries are
nonnegative, and the complement's positive mass cannot
increase, so the plane loss is at most the positive mass loss.

For a complement weight \(\kappa\), denote the positive and negative masses by \(m^+=\sum_i(\kappa_i)_+\) and \(m^-=\sum_i(-\kappa_i)_+\). At a vector step the change in the sum of the entries is 
\[
 \Delta\!\sum_i\kappa_i=\Delta m^+-\Delta m^-,
\]
 the positive mass gain plus the negative mass loss. Interlacing bounds its
positive mass gain by the last plane entry before the step.
Reading that entry backwards through the last \(r-1\) vector
steps bounds it by the first entry just after \(u\). Geometric growth of the vector degrees gives
\[
 \text{first entry just after }u\ \le\!\sum_{x\le u,\ x\text{ a vector}} B_x\le 2B_u.
\]
Using the bound
\[
 2B_u\le4a/P_{j,u},
\]
the complement
gain is at most
\[
       \text{negative mass loss}+2B_u
                      \le \Xi_x+4a/P_{j,u}.
\]

\proofstep{The loss from the mass test.}
For a violation of the mass condition let \(U\) be the sum
of preceding degrees of the opposite sign. If there is any
mass of that sign, at least one such pre-outgoing degree has
appeared. Its largest degree is at least \(k\) and at most
\(4\beta_x/p\). The geometric sum gives
\[
 U\le5\beta_x/p.
\]

The total mass of the sign before this step is at most \(U\)
by the deficit inequality at \(T\ge S\). Hence for \(T\ge U\)
the deficit of that sign grows by more than \(\Xi_x\). This
includes negative mass losses at outgoing steps, since removing
an entry equal to \(S\) preserves the negative mass.

\proofstep{The loss from the last-entry test.}
For a violation of the last-plane-entry condition, put \(B=B_u\).
If \(t\) outgoing steps have passed at this cut, \(u\) is the
oldest of the
\[
 h=r-t
\]
most recent pre-outgoing vectors.
The sum of positive rows after the first \(h\) in the tail is
at most the sum of all vector degrees strictly before \(u\).
At each of those \(h\) vector additions, interlacing bounds the
positive sum below row \(i+1\) after the addition by the sum
below row \(i\) just before. Covectors and the later outgoing
trims can only decrease these sums. Thus the positive sum after
row \(h\) is at most \(B/(\rho-1)\). Row \(h\) has positive part
less than \(2B/3\).

For
\[
 B\le T\le13B,
\]
the \(h\) latest pre-outgoing
vectors contribute at least \((h-1)T+B\) to the reference sum.
In the clipped tail, the first \(h-1\) rows contribute at
most \((h-1)T\), row \(h\) at most \(2B/3\), and the rest at
most \(B/(\rho-1)\). Subtracting gives
\[
 \delta_+(T)\ge \big((h-1)T+B\big)
                 -\big((h-1)T+2B/3+B/(\rho-1)\big)\ge B/4
\]
by the stipulated growth of the vector degrees.

\proofstep{Converting exceptional events into lost weight.}
To convert these deficits into a weight loss, set
\(J_k(T)=\sum_{l=1}^{k-1}(T+l)^{-2}\). For \(b\ge0\),
\[
 \log D_k(b)=\int_0^\infty\min(b,T)J_k(T)\,dT,
\]
where the left side for real \(b\) denotes the logarithm of the
binomial product. The equality follows by integrating each
\((T+l)^{-2}\) separately.

At the terminal complement weight
\(\lambda\) in dimension \(K\), the positive and negative lists
bound its dimension by
\[
 \dim V_\lambda\le\prod_h D_k(\alpha_h^+)D_k(\alpha_h^-).
\]
Indeed each Weyl factor is bounded by
\[
 1+\frac{\lambda_i-\lambda_j}{j-i}
 \le\left(1+\frac{\max(\lambda_i,0)}{j-i}\right)
     \left(1+\frac{\max(-\lambda_j,0)}{j-i}\right).
\]
Multiplying gives the bound with \(D_K\), which is at most
\(D_k\) for each nonnegative argument. Therefore
\(\log(R/\dim V_\lambda)\) is at least the integral of
\((\delta_++\delta_-)J_k\) at the last cut. If the mass bound
fails, this integral is at least
\[
 \Xi_x\sum_{l=1}^{k-1}\frac1{U+l}
 \ge \frac{\mathcal L(k-1)}{9k}\ge\frac{\mathcal L}{16},
\]
using \(k\le4\beta_x/p\). If the plane cutoff fails, it is at
least
\[
 \frac B4\sum_{l=1}^{k-1}
     \left(\frac1{B+l}-\frac1{13B+l}\right)
 \ge\frac{11(k-1)}{104}\ge\frac{k}{16}.
\]
Here \(B\ge k\), and both estimates are independent of the
remaining path and its strip degrees.

At each cut the two lists of nonzero parts have lengths at
most \(m\), since each pre-outgoing strip introduces at most
one new part of its sign. The outgoing interlacing keeps that
bound after removing the heads. Each part is an integer in
\([0,S]\). The lists padded to length \(m\), and the time of
the cut, determine the whole shape. There are at most
\((1+S)^{2m^2}\) shape paths because the zero start is fixed.
Thus the exceptional weight is at most
\[
 R\exp\bigl(2m^2\log(1+2aN_s)-\min(k,\mathcal L)/16\bigr)
 \le R a^{-N_s^2}.
\]
The last inequality uses \(1+2aN_s\le a^2\), supplied by
\eqref{eq:plane-count-size}, and the bounds on \(k,\mathcal L\).
This proves \eqref{eq:plane-good-path-count} for each choice
of intervening digits with the same constants.
\end{proof}


\paragraph{Sum over outgoing rows and digits.}
The last gap before the first outgoing position exceeds \(L_1\).
Thus each pre-outgoing \(x\) has
\(P_{j,x}\ge p^L p^{L_1+1}/H_j\). The positions are distinct, and
\(B_x\le 2a/P_{j,x}\). Summing the geometric series and using
\(a/p^L<8pYk^2N_s\) gives
\[
 S\le \frac{4aH_j}{p^L p^{L_1+1}},\qquad
 4YS\le H_j .
\]

For each outgoing row of a fixed path, the congruence
\(d_i+\sum_h v_{i,h}\equiv n_i\pmod{H_j}\) has exactly
\(H_j^{r-1}\) solutions with \(0\le v_{i,h}<H_j\). Each coordinate
is uniform among these solutions since \(r\ge2\). The union
bound removes at most \(2rS H_j^{r-2}\) solutions on requiring
\(v_{i,h}<H_j-2S\) for all \(h\). Each surviving solution has
\(d_i+\sum_h v_{i,h}<rH_j\), and so supplies a permitted outgoing
degree. The count per row is at least
\(H_j^{r-1}(1-r/(2Y))\) when this factor is nonnegative.

\begin{proposition}[Rank after summing the coarsened families]\label{prop:coarse-plane-rank}
With the admissible frame degrees and path count
\eqref{eq:plane-good-path-count}, the coarsened generators
span a residue space \(\overline\Lambda\) with
\[
 \dim\overline\Lambda\ge q^{-rK}(1-Nr^2 k/Y)\prod_xD_k(b_x).
\]

\end{proposition}

\begin{proof}
\proofstep{Summation using the degree at the outgoing row.}
For fixed digits the number of decorated words depends only on the outgoing degree
of its own segment. The incoming padding changes with adjacent
outgoing degrees, but the independence and membership just
proved hold at every such degree. Hence the sum of products
of word counts over the outgoing choices factors across segments.
For the digit sum at a top between frames, dividing each of the
\(k\) exponents of a monomial of degree \(a'_x\) by \(P_{j,x}\)
gives the identity \(\sum_{\ell_x} D_k(n_x)=D_k(a'_x)\).

At a frame top there are exactly \(P_{j,x}^{k-1}\) digit choices
by the sum congruence. Its incoming path factor in the next
segment is \(D_k(W_{j+1,x}-m_0)\), so its full contribution is
at least
\[
 \frac{1-r/(2Y)}{H_j^K}\,
 (P_{j,x}H_j)^{k-1}D_k(W_{j+1,x}-m_0).
\]
This is where the complement dimension \(K\) in the exponent
arises. The product of the factors \(H_j^{-K}\) over all \(r\)
rows in each frame is \(q^{-rK}\).

For the other factors, \(E'_x+E_x\le5b_x/(kY)\).
The binomial product formula and Bernoulli's inequality give
\(D_k(a'_x)\ge(1-5/Y)D_k(b_x)\). For an incoming radix
\(P'=P_{j,x}H_j\), put \(B'=W_{j+1,x}-m_0\). The floor and
padding lose at most \(4kN_s\) in the coarse degree.

Since
\(a'_x/P'\ge 4Yk^2N_s\),
\[
 P'B'\ge a'_x-4kN_sP'\ge(1-1/(kY))a'_x,\qquad
 (P')^{k-1}D_k(B')\ge(1-1/Y)D_k(a'_x).
\]
For the latter comparison use \(P'(B'+i)\ge(1-1/(kY))(a'_x+i)\)
for each \(1\le i\le k-1\). If \(1-Nr^2k/Y\le0\), the
claimed lower bound is automatic. Otherwise every loss factor
is in \([0,1]\). Since \(a\ge32Y\), \(J\le N\), and \(rJ\le N\),
the product of all loss factors is at least
\[
 1-N(5+1+r/2+1/32)/Y\ge 1-Nr^2k/Y.
\]

The next bound counts independent residue symbols in the
coarse radices. The following paragraph then checks that the
same families are admitted in the original radices, as needed
for Lemma~\ref{lem:lift-plane-symbol}.

To see that the dimension counts add, recover the choices from
the exponent at each top. For a coordinate \(h\) in a shared
frame row, write its true exponent as
\[
 \ell_{x,h}+P_{j,x}\bigl(a_{x,h}+H_j b_{x,h}\bigr).
\]
The digit \(\ell_{x,h}\) is the remainder modulo \(P_{j,x}\).
The outgoing exponent \(a_{x,h}\) is less than \(H_j\)
by the coordinate condition, so division of the quotient by
\(H_j\) recovers \(a_{x,h}\) and \(b_{x,h}\). In a complement
coordinate \(b_{x,h}=0\). The sum of the outgoing exponents
recovers \(e_x\). At unshared tops there is no second quotient.
Hence multiplication induces an injective map on the monomial
tuples, and distinct choices of digit and outgoing degree give
disjoint monomial supports. By
\eqref{eq:plane-good-path-count} and the count just proved,
\begin{equation}\label{eq:plane-coarse-rank}
 \dim \overline\Lambda\ \ge\ q^{-rK}(1-Nr^2k/Y)\prod_xD_k(b_x).
\end{equation}

\end{proof}

\paragraph{Return to the original radices.}
\label{par:plane-rank-transfer}

The count so far used coarsened digit classes. Returning to the original radices determines the rank available for the actual symbols.

\begin{proposition}[Rank of symbols at the original radices]\label{prop:plane-original-symbol-rank}
In the setting of Proposition~\ref{prop:coarse-plane-rank},
uncoarsening and multiplication by the reserved determinants
give integral original symbols of degrees \(b_x-E'_x\)
satisfying the required prefix tests. Their span contains a direct summand of rank \(r_{\rm orig}\) satisfying the bound in \eqref{eq:plane-coarse-rank},
\[
 r_{\rm orig}\ge q^{-rK}(1-Nr^2k/Y)\prod_x D_k(b_x).
\]

\end{proposition}

\begin{proof}
It suffices to return the products of the counted words to
the original symbol module. Use \(P^0_{j,x}=p^{z_j-x}\),
\(F_j=P_{j,x}/P^0_{j,x}=p^{R_j}\). Euclidean division in the unique late segment gives
\[
 \ell_{x,h}=\ell^0_{x,h}+P^0_{j,x}u_{j,x,h},\qquad
 0\le\ell^0_{x,h}<P^0_{j,x},\quad 0\le u_{j,x,h}<F_j .
\]

To
uncoarsen a counted segment product, raise its plane and
complement factors to \(F_j\). Multiply each factor by the
corresponding coordinate monomial of exponents \(u_{j,x,h}\)
in the late rows. Multiply the plane factor also by the
ordinary determinants of that segment with powers \(g_{j,x}\).
All coefficients are integers. Indeed, for an integral factor \(A\) with integer coefficients, Frobenius gives
\[
 \overline{A(X)^{F_j}}=\overline{A}(X^{F_j})
\]
in the residue field, with coordinatewise powers on the right. Applying the original radices in every
segment and multiplying by the remaining digits \(\ell^0\)
therefore gives \(\overline{\mathcal D}\) times the coarsened
product. In the integral polynomial the total degrees are
the coarsened true degrees plus the exact determinant degrees
\(E_x\), hence equal to \(b_x-E'_x\).

\proofstep{The inequalities for the restored products.}
We check the original inequalities on these integral products.
The plane and complement letter supports are unchanged, since
the digits are placed only in late rows and the determinants
use earlier plane vectors. There are at most \(kN_s\) coordinate digits in a prefix and each has size at most \(F_j-1\). Their combined adverse change in separation is therefore at most
\[
 \sum_{x,h}u_{j,x,h}\le kN_s(F_j-1).
\]
An odd plane coordinate can lower the plane
minimum by at most its exponent, and an even complement
coordinate can raise the complement maximum by at most its
exponent.

The other coordinate types do not decrease the
separation. The determinant associated to a given cut lies
wholly before that cut, so it contributes \(g_{j,x}\) to the
minimum. All other determinants have nonnegative minimum
since they contain only vectors. Adding the determinant gain to the worst digit loss gives, at this cut,
\[
 \text{separation}\ge F_jm_0-kN_s(F_j-1)+g_{j,x}>m_*+g_{j,x}.
\]
For the strict inequality use the separation margin in \(m_0\) and
\(F_j>a/(8pYk^2N_sQ)>m_*\) from \eqref{eq:plane-count-size} and
the definition of \(L\).

\proofstep{The rank lower bound at a plane.}
For the absolute plane lower bound put \(M=b_u/P_{j,u}\) at
the earliest preceding vector \(u\). The count assumptions
include \(k\ge C_kN_s^2\log a\) with \(C_k\ge1088\). In
particular \(a'_u\ge9b_u/10\) by the bound on \(E_u+E'_u\).
Since \(u\) is pre-outgoing, the degree \(B_u\) used by the
count satisfies \(B_u\ge 9M/10-4kN_s\) for either an
incoming or an intervening row. Also \(M\ge8Yk^2N_s\ge
40kN_s\). The plane minimum before the cut in the uncoarsened
word is at least
\[
 F_j(2B_u/3-kN_s)\ge F_j M/2=\frac{b_u}{2P^0_{j,u}}.
\]

For the normal cutoff at \(x\), the extra digit degree is at
most \(kF_j\). The gap \(x-u>L_1\) and \(F_jP^0_{j,u}\le A\)
give \(kF_j\le a/(p^{L_1}P^0_{j,x})\). The term at an even
\(x\) from the coarsened normal bound is
\[
 4a/P^0_{j,u}\le a/(p^{L_1}P^0_{j,x}).
\]
Thus the total
above \(\mathcal L a/(pkP^0_{j,x})\) is bounded by
\(a/(p^{L_1}P^0_{j,x})\) at an odd top and by twice that at
an even top. These imply both of the original order bounds.

The original symbols produced this way have the thresholds,
normal cutoffs and degrees needed by the precision induction.
After the original radix substitutions, the total normal
degree at each top is bounded by
\(\mathcal L a/(pk)+10a p^{-L_1}+kQ\). Indeed there is one
late occurrence, the incoming contribution is tangential,
and the final digit costs at most \(kQ\). The same upper bound
holds on supports of sums, so the total over all tops is at
most
\[
 N(\mathcal L a/(pk)+10a p^{-L_1}+kQ)
\]
on the whole
retained lattice.

The determinant product has nonzero reduction. For each of
its factors the preceding \(r\) even positions are distinct
modulo \(n'\). An interval of length \(n'\) starting immediately
after the earliest of those positions contains all \(r\) even
positions of any given frame modulo \(n'\), each once. If the
last position were at least that far from the earliest, there
would be at least \(r\) later preceding positions in the same
segment. Therefore the determinant monomials remain distinct
under the positive-power substitutions and have coefficients
\(\pm1\).

Since the residue polynomial ring is a domain, multiplication by the nonzero determinant product has zero kernel,
\[
 \overline{\mathcal D}\sum_i a_iW_i=0\quad\Longrightarrow\quad\sum_i a_iW_i=0.
\]
Hence the
products counted in \eqref{eq:plane-coarse-rank} remain independent
after multiplying by \(\overline{\mathcal D}\). Complete their
independent reductions to a basis of the ambient fixed-degree
residue space. Lifting the additional basis vectors to the
integral monomial lattice gives a basis by a change of basis
with unit determinant. This gives an integral direct summand
of the indicated rank generated by original symbols.
\end{proof}
